\section{Problem Formulation}

\paragraph{Assumptions}
For Variational Automation tasks that will be repeated over extended periods, we assume the workcell environment, robot, and sensors are known and fixed. Furthermore, we assume the range of potential objects and the range of initial object poses are known. These assumptions are part of the VA setting rather than oracle information: unlike generalist robotics, VA tasks are defined by a known workcell and bounded operating envelope, allowing systems to use object models, calibrated sensors, and reusable automation skills when available.

\paragraph{Variational Automation Task Class}
We formalize a Variational Automation (VA) Task using the tuple $\mathcal{T} = \langle \mathcal{L}, \mathcal{E}, \mathcal{R}, \mathcal{O}, \mathcal{X}, \mathcal{B}, \mathcal{J} \rangle$:

\begin{itemize} [nosep, leftmargin=*]
    \item $\mathcal{L}$, Language Space: Natural language instructions and semantic descriptors that describe desired task behavior, providing high-level context for decomposing the task into segments.
    \item $\mathcal{E}$, is a known stationary workspace environment (e.g., workcell) that establishes world frame $\mathcal{W}$ and occupancy map $\mathcal{M}_E$, providing the support surfaces that constrain the feasible pose space of all entities to the non-colliding subset of $SE(3)$.
    \item $\mathcal{R}$, is the {robot and sensor configuration}, specifying the robot URDF, joint limits, gripper configuration, camera spec, and camera placement.
    \item $\mathcal{O}$, Object Set: The set of all possible objects in the task, comprising both rigid objects $\{o_i\}$ with 3D models $\mathcal{M}_i$ and articulated entities $\{\kappa_j\}$ (e.g., drawers, knobs) with defined kinematic joint limits $[\theta_{\min}, \theta_{\max}]$.
    \item {$\mathcal{X}$, State Space:}  The product space of the robot joint configurations, the $SE(3)$ poses of all rigid objects, and the kinematic states of articulated entities: $\mathcal{X} = \mathcal{C}_{robot} \times SE(3)^n \times \mathbb{R}^m$.
    \item {$\mathcal{B}$, Belief Space:} Describes the distribution of objects and poses in instances of the task. A specific instance is sampled $\mathbf{x}_i \sim p(\mathbf{x} \mid\mathcal{X})$, introducing variations such as: (i) \textit{Structured Priors:} Position sampled uniformly over a volume $\mathcal{V}$ ($\mathbf{x}_i \sim \text{Uniform}(\mathcal{V})$) and orientation ranges, and (ii) \textit{Empirical Distributions:} Multi-modal distributions estimated from real-world demonstrations or perception (e.g., point cloud registration).
    \item $\mathcal{J}$, Reward function: A multi-objective reward function used to evaluate success and efficiency, defined as:$\mathcal{J} = w_s \cdot \mathbb{I}(\text{success}) + w_t \cdot \Phi$ where $\mathbb{I}(\cdot)$ is the success indicator, $\Phi$ is the throughput (success rate / cycle time), and $w_s, w_t$ are weighting factors.

\end{itemize}

Given a task $\mathcal{T}$, a task instance $\tau_i = \langle o_i, x_i \rangle$ is drawn from $\mathcal{O}$ and $\mathcal{B}$ respectively.
   
\paragraph{Policy Representation via Directed Execution Graphs}

Graph-as-policy (GaP) represents a robot policy $\pi(a \mid \mathbf{x}, \mathcal{T})$  as a directed computation graph $\mathcal{G} = (V, E)$ to complete all task instances   $\forall \tau_i = \langle o_i \subseteq \mathcal{O}, x_i \sim \mathcal{B} \rangle \in \mathcal{T}$. The graph $\mathcal{G}$ consists of modular, atomic functional units called \emph{nodes}, connected by edges that dictate both data flow and execution logic. The components are defined as follows: \textbf{Nodes ($V$):} Each node $n \in V$ represents a functional primitive for manipulation, perception, or a segment of executable code. Each node encapsulates a single, well-defined robotic operation with a typed input/output signature.
Nodes can be grouped into \textbf{skills},  a natural-language specification that instructs an LLM agent how to configure and compose a set of atomic nodes on the execution graph for a defined sub-task. 
\textbf{Edges ($E$):} A directed edge $e = (n_i, n_j) \in E$ represents a data or logic dependency. {Data edges} route a producer node's output to a consumer node's input (e.g., feeding the output of an object-centering node into a planar-surface-orienting node) and implicitly induce execution order through their dependencies; independent branches may execute concurrently. {Control edges} are condition branches and carry a predicate over node outputs.

% An individual task instance is a realization sampled from the belief space $\mathcal{B}$, characterized by a specific instruction $l \in \mathcal{L}$ and a scene configuration $\mathbf{x} \sim p(\mathbf{x} \mid \mathcal{T})$.

% \paragraph{Task Instance and Belief Space}

% The belief space $\mathcal{B}$ describes the distribution over possible task instances  for a given task class. A specific task instance is sampled from $\mathcal{B}$, introducing variations that the agent must perceive or be robust to.
% For each rigid object $o_i \in \mathcal{O}$, its state is defined by its pose $\mathbf{x}_i \in SE(3)$. A task instance defines a distribution over these poses: $\mathbf{x}_i \sim p(\mathbf{x} \mid \mathcal{T})$. This distribution can take several forms:(a) \textbf{Structured Priors:} For position perturbation, the object center may be sampled uniformly over a volume $\mathcal{V}$ such that $\mathbf{t}_i \sim \text{Uniform}(\mathcal{V})$, with orientation $\theta_i$ sampled from a specified range. (b) \textbf{Empirical Distributions:} Distributions estimated from real-world demonstrations or perception outputs (e.g., multi-modal distributions from point cloud registration). 



% \paragraph{Augmented Uncertainty} 
% The belief space $\mathcal{B}$ is not limited to rigid poses; it may also capture uncertainty in articulated joint states $\mathbf{q}_j$, object identity, semantic segmentation, or system noise including camera calibration errors and depth estimation variance. Formally, the agent must operate over the joint distribution of all stochastic elements within the workspace.





\paragraph{Problem Definition}

Given a task class $\mathcal{T} = \langle \mathcal{L}, \mathcal{E}, \mathcal{R}, \mathcal{O}, \mathcal{X}, \mathcal{B}, \mathcal{J} \rangle$, GaP must synthesize a robust execution graph $\mathcal{G}^*$ that generalizes across the entire belief space $\mathcal{B}$. Formally, for any task instance $\tau_i \in \mathcal{T}$ and given a real-time scene observation $\mathcal{I}$ (comprising multi-view image sets from static and wrist cameras), the graph executor interprets
the computation graph $\mathcal{G}^*$ by invoking its skill nodes and following its data and control edges.  This execution induces a closed-loop robot policy $\pi_G(a \mid \mathcal{I})$, which maps observations and graph state to robot actions in order to achieve the goal state specified by $\mathcal{L}$.  We define the optimized graph $\mathcal{G}^*$ as one that maximizes performance across all instances in the variational class:~$\mathcal{G}^* = \arg\max_{\mathcal{G}} \mathbb{E}_{x_i \sim \mathcal{B}} \left[ \mathcal{J}(\pi(a\mid\mathcal{I}, \mathcal{G})) \right]$. 

The optimized graph is then sent to an edge-based interpreter for repeated execution on the robot.

% where $\mathcal{J}(\pi(a\mid\mathcal{I}, \mathcal{G}))$ is the {performance objective function}. In our framework, $\mathcal{J}$ is defined as a multi-objective return: $\mathcal{J} = w_s \cdot \mathbb{I}(\text{success}) + w_e \cdot \Phi$,
% where $\mathbb{I}(\cdot)$ is the indicator function for task completion, $\Phi$ is the throughput (success rate / cycle time), and $w_s, w_e$ are weighting factors. This objective ensures the synthesized graph $\mathcal{G}^*_\mathcal{T}$ is not only successful across the distribution of instances $x_i \sim \mathcal{B}$ but also minimizes execution time.

% Given scene observation $\mathcal{I}$ comprised of multi-view image sets (e.g., static side-view and eye-in-hand wrist cameras), the goal is to synthesize a robot policy $\pi(a \mid \mathcal{I}, \mathcal{T})$ that maps the visual input and task specification $\mathcal{T}$ to an action sequence. We evaluate the performance of the policy using two primary metrics: (i) Success Rate ($r_{\tiny \mbox{success}}$): The ratio of successfully completed trials to total attempts, and  (ii) Throughput: Defined as $\frac{r_{\tiny \mbox{success}}}{t_{\tiny \mbox{cycle}}}$, where $t_{\tiny \mbox{cycle}}$ is the average time taken per task execution, measuring the efficiency of the learned behavior.


% where a detection and classification pipeline identifies the relevant environment $\mathcal{E}$ and object sets $\{\mathcal{O}, \mathcal{K}\}$ from an initial observation. 
% While $\mathcal{T}$ is known, the specific configuration of the task instance within the belief space $\mathcal{B}$ must still be resolved by the agent.
% \eric{metric, success rate / cycle time = throuhgput}

% \begin{figure}
%     \centering
%     \footnotesize
%     \includegraphics[width=\linewidth]{figures/main.png}
%     \caption{\textbf{GaP} is a multi-agent harness framework. Given a natural language task specification, it utilizes a collaborative multi-LLM architecture and a Modular Open Robot Skill Library (MORSL) to autonomously decompose and synthesize an interpretable execution  graph. It uses simulation-based self-learning by automatically instantiating parameterized internal physics environments to rehearse the task motion, and uses closed-loop state and contact feedback to iteratively refine the graph structure and node parameters. }
%     \label{fig:main figure}
% \end{figure}


\section{Graph-as-Policy}
\label{sec:design}
The GaP architecture is illustrated in Figure \ref{fig:scene-title}.  As an example, consider the Make Popcorn VA task. Given a natural language specification and geometric object models, the multi-agent GaP harness first partitions the high-level objective into semantic segments such as ‘turn on the knob’, ‘pick up the popcorn pan’, etc. Next, drawing from the Modular Open Robot Skill Library (MORSL), GaP maps these segments into atomic robotic skills to synthesize an initial computation graph $\mathcal{G}$.  During self-learning, the computation graph is evaluated using an internal simulation to iteratively refine the execution graph topology and node parameters prior to deployment. The optimized robot computation graph is then sent to an external interpreter for repeated execution on the physical robot.

%\subsection{Behavior Graph Generator: Multi-Agent LLM Harness }

%To autonomously synthesize the policy graph from a natural language instruction, GaP uses a hierarchical, multi-agent architecture to generate robot behavior based on skill availability.
%Specifically, the GaP framework consists of three specialized components: a \textit{Behavior Agent}, multiple \textit{Skill Agents}, and a \textit{Verification Agent}.
% a \textit{Behavior Agent} that generates a skill-aware task plan, a \textit{Skill Agent} that translates specific sub-task specifications into functional local subgraphs, and a \textit{Verification Agent} that generates the postconditions for the nodes and skills to validate the semantic correctness and physical feasibility for Sec. \ref{sec:optimization}.
% The generation process is orchestrated by a  Behavior Agent. 
%Given a high-level task specification, the Behavior Agent acts as the primary planner, decomposing the global objective into a sequence of logical sub-tasks based on the available skills. This decomposition yields a high-level Behavior Graph ($\mathcal{G}_B$), which captures the semantic dependencies and control flow of the task prior to parameterizing the low-level execution details.
%For each identified sub-task in $\mathcal{G}_B$, the Behavior Agent spawns a dedicated Skill Agent, which receives the Behavior Agent's localized instruction alongside the corresponding skill specification. Acting as a domain expert, the Skill Agent generates the necessary programmatic logic to instantiate and connect the appropriate atomic nodes, outputting a functional, localized subgraph.
%After Skill Agents generate their respective subgraphs, the Behavior Agent aggregates and wires them together to form the complete execution graph $\mathcal{G}$. During this aggregation phase, the Behavior Agent performs a static type-checking and structural verification pass across all edge interfaces to guarantee that all required nodes are available and that the system's cross-component connectivity is structurally sound in terms of strict input-output type matching. Additionally, the Verification Agent generates sub-task post-conditions for graph nodes and skills to be used in validation for Sec. \ref{sec:optimization}.



% spanning the full perception, planning and control spectrum, such as retrieving the latest RGB-D camera frame, running inference with a learned foundation model (e.g., a Vision-Language-Action model or dedicated object detector), moving a manipulator to a target Cartesian pose, or actuating a gripper. 


\subsection{Modular Open Robot Skill Library (MORSL)}
The MORSL library uses agentic tool-use conventions (e.g., Anthropic's
\texttt{Skill.md}~\cite{anthropicskill}) with extensions for graph declarations. Each skill declares its inputs, outputs, semantic parameters, and pre-conditions, so the agent can decide both when to invoke it and how to wire it into the graph. Examples of the 51 initial MORSL skills include perception (SAM2~\cite{ravi2024sam}/3~\cite{carion2025sam3segmentconcepts}, Grounding DINO~\cite{liu2023grounding}, OWL-ViT~\cite{minderer2022simple}, Molmo~\cite{deitke2024molmo}, general-purpose VLM~\cite{geminiroboticsteam2025geminiroboticsbringingai, kamath2025gemma}; 15 skills), grasp planning (Contact GraspNet~\cite{sundermeyer2021contact}, GraspGen~\cite{murali2025graspgen}, M2T2~\cite{yuan2023m2t2}; 5 skills), motion planning (cuRobo~\cite{sundaralingam2023curobo} and cuRobov2~\cite{sundaralingam2026curobov2}; 8 skills), 2D and 3D vision  utilities with NumPy and OpenCV (e.g., DBSCAN for point cloud processing; 15 skills), and 8 additional verification and control primitives including
(*) {Cartesian Linear Motion Planning with CuRobo}; 
(*) {Robot Operating System (ROS) Translator };  
(*) {Visuomotor Interactive Perception Policies}.  The full catalog can be found in the Appendix A.


% PyRoKi front-loads cheap feasibility checks (reachability, null-space posture) to prune infeasible
% grasps before invoking expensive planning; cuRobo then produces collision-free trajectories against the current scene. Exposing both as separate skills lets the agent fall back gracefully on planning failure—relaxing tolerances, reseeding, or requesting
% re-perception—rather than treating motion as an opaque "move to pose" call.




\subsection{Self-Learning through Internal Simulation Rehearsal}
\label{sec:optimization}
% \subsection{Real-to-Sim Scene Specification}
% LLM agent design, recgen, etc
\renewcommand{\algorithmcfname}{Self-Learning-Algorithm} % algorithm2e's caption name (\floatname is from the float package, no longer loaded)
% \subsection{Skill-aware graph update with Physical feedback}
\begin{wrapfigure}{r}{0.55\textwidth}
\vspace{-24pt}
\begin{minipage}{0.55\textwidth}
\begin{algorithm}[H]
\scriptsize 
\DontPrintSemicolon
\caption{Rehearsal-based Graph Optimization}
\label{alg:graph_opt}
\SetKwInOut{KwInput}{In}
\SetKwInOut{KwOutput}{Out}
\noindent\textbf{Input:} $\mathcal{T} = \langle \mathcal{L}, \mathcal{E}, \mathcal{R}, \mathcal{O}, \mathcal{X}, \mathcal{B}, \mathcal{J} \rangle$,  Iterations $M$, Parallel Rollouts $N$
\noindent\textbf{Output:} Optimized Task Graph $\mathcal{G}^*$
\BlankLine
\textbf{Initialization:}\;
$\mathcal{G}_0 \leftarrow \mathrm{Graph\_Init}(\mathcal{L})$\;
\quad\tcp*[r]{\tiny Initial graph from language}
$\mathcal{S} \leftarrow \mathrm{Build\_Scene}(\mathcal{G}, \{\mathcal{G}_i\})$\;
\quad\tcp*[r]{\tiny Instantiate sim environment}
\For{$j \leftarrow 1$ \KwTo $M$}{
    \tcp{Step 1: Scene Variational Sampling}
    $\{\hat{s}_i\}_{i=1}^N \sim \mathcal{B}$\;
    \quad\tcp*[r]{\tiny Sample $N$ instances}
    \tcp{Step 2: Parallel Rehearsal}
    \ParallelFor{$i \leftarrow 1$ \KwTo $N$}{
        $\tau_i \leftarrow \mathrm{Rehearsal}(\hat{s}_i, \mathcal{G}_{j-1})$\;
        \quad\tcp*[r]{\tiny Rollout policy}
        $F_i \leftarrow \mathrm{Analyze_\_Failure}(\tau_i, \mathcal{G}, \{\mathcal{G}_i\})$\;
    
    }
    \tcp{Step 3: Graph Refinement}
    $\mathcal{G}_j \leftarrow \mathrm{Graph\_Update}(\{F_1, \dots, F_N\})$\;
    \quad\tcp*[r]{\tiny Optimize via LLM}
}
\Return $\mathcal{G}^* \leftarrow \mathcal{G}_M$\;
\end{algorithm}
\end{minipage}
\vspace{-8pt}
\end{wrapfigure}
To iteratively learn to improve a graph using task instances, we execute the computation graph, provide feedback, update the graph, and iterate. 
As described in self-learning Algorithm~\ref{alg:graph_opt}, GaP uses Isaac simulator~\cite{NVIDIA_Isaac_Sim} as an internal simulation to render the visualization for the perception nodes, simulate physics, and compute contacts. 
To generate feedback, GaP registers the robot and object state before and after each rehearsing node to compute differences and infer the motion outcome. 
The optimization process begins with task instance sampling, where the system generates $N$ parallel task instances  $\{{\tau}_i\}_{i=1}^N$. These task instances are instantiated by sampling from the belief space $\mathcal{B}$, which represents the probability distribution over potential object poses, initial states, and kinematic configurations within the workspace $\mathcal{E}$. 
GaP generates the initial graph $\mathcal{G}_0$ through multi-agent harness. 
Then, a parallel rehearsal starts across these $N$ scenes to evaluate graph performance. In instances where the rollout fails to meet success criteria, GaP analyzes the physical execution data to isolate geometric root causes within the occupancy or articulated entities. The graph is updated by the agents with Graph Update that iteratively modifies the graph architecture (e.g., swapping functionally equivalent nodes, changing edges, and updating code parameters) until the system performance reaches a plateau or terminates with failure reasons.


% \begin{algorithm}[H]
% \caption{Skill-Aware Graph Update with Physical Feedback}
% \label{alg:gap_update}
% \begin{algorithmic}[1]
% \REQUIRE Task Class $\mathcal{T}$, Initial Graph $\mathcal{G}_0$, Number of Scenes $M$, Max Iterations $N$, Success Threshold $\gamma$, Scene Generation Agent $\mathcal{A}_{scene}$, Graph Update Agent $\mathcal{A}_{update}$
% \ENSURE Final Executable Graph $\mathcal{G}^*$

% \STATE $\{S_m\}_{m=1}^M \leftarrow \mathcal{A}_{scene}(\mathcal{T}, M)$
    
% \FOR{$i = 0$ \TO $N-1$}
%     \STATE $\mathcal{O}_i \leftarrow \text{Rollout}(G_i, \{S_m\}_{m=1}^M)$
    
%     \IF{$\text{SuccessRate}(\mathcal{O}_i) \geq \gamma$}
%         \RETURN $\mathcal{G}_i$
%     \ENDIF
    
%     \STATE $f \leftarrow \text{AnalyzeFailure}(\mathcal{O}_i)$
%     \STATE $\mathcal{G}_{i+1} \leftarrow \mathcal{A}_{update}(G_i, f)$
% \ENDFOR

% \RETURN $\mathcal{G}_N$
% \end{algorithmic}
% \end{algorithm}
